\chapter{Render output golden tests}
\label{ch:golden}

\begin{aside}
The diff is the contract: same input, same bytes on the wire.
A golden test pins that contract.
\end{aside}

\section{Golden vs snapshot}

A snapshot test (see Chapter~\ref{ch:snapshot}) checks the rendered
canvas. A golden test checks the emitted ANSI byte stream.

Two layers, two tests:

\begin{itemize}[leftmargin=*]
  \item Snapshot catches ``the screen looks wrong''.
  \item Golden catches ``the wire bytes changed''.
\end{itemize}

A refactor of the emit step that produces equivalent-but-
different ANSI passes snapshot, fails golden. Whether that
fail is intentional is a judgement call.

\section{Capturing a golden}

\begin{lstlisting}
TEST_CASE("frame emits expected ANSI") {
    StringBackend back;
    Renderer r(back);
    auto el = my_panel();
    r.render(el);
    EXPECT_BYTES_EQ(back.output, "fixtures/panel.ansi");
}
\end{lstlisting}

The fixture is a raw ANSI dump.

\section{Diffing ANSI}

A diff between ANSI streams is harder to read than a diff
between cell grids. \maya{}'s golden helper pretty-prints
escape sequences:

\begin{lstlisting}
ESC[?1049h ESC[?25l ESC[H ESC[31mhello ESC[0m
\end{lstlisting}

becomes:

\begin{lstlisting}
ENTER_ALT_SCREEN
HIDE_CURSOR
CURSOR_HOME
SGR(fg=red)
text "hello"
SGR_RESET
\end{lstlisting}

Reviewing this in a PR is feasible.

\section{Volatile bytes}

Some bytes vary across runs:

\begin{itemize}[leftmargin=*]
  \item Cursor position based on terminal size detection.
  \item Time-based escape sequences (very rare; none in
    \maya{}).
  \item Performance-related ANSI like DEC 2026 only
    when supported.
\end{itemize}

Goldens use a fixed backend with fixed caps and fixed size
to avoid this.

\section{What a golden test catches}

\begin{itemize}[leftmargin=*]
  \item Accidental SGR re-emission (style cache bug).
  \item Wrong escape syntax (missing semicolon, off
    parameters).
  \item Order changes (e.g., cursor move after SGR vs
    before).
  \item Whitespace handling.
\end{itemize}

These are exactly the bugs that survive a snapshot test ---
the cells look right but the wire is wrong.

\section{Maintaining goldens}

Treat \code{.ansi} files as code. Review changes. When a
deliberate optimisation changes bytes, update intentionally
with notes in the commit message.

\section{Recap}

\begin{recap}
\begin{itemize}[leftmargin=*]
  \item Golden test pins the wire output.
  \item Pretty-print ANSI for readable diffs.
  \item Fix backend caps and size to keep goldens stable.
  \item Catches bugs snapshots miss.
\end{itemize}
\end{recap}

\section*{Workshop}

\begin{enumerate}[leftmargin=*]
  \item Capture a golden for your app's first frame.
  \item Make a perf change that emits the same cells in
    different byte order. Update golden.
  \item Verify the pretty-printer correctly decodes every
    escape your app emits.
\end{enumerate}
